\documentclass[12pt,a4paper]{article}

\usepackage{indentfirst}
\usepackage[utf8]{inputenc}
\usepackage[T1]{fontenc}
\usepackage[brazil]{babel}
\usepackage{amsmath}
\usepackage{geometry}
\usepackage{setspace}
\usepackage{listings}
\usepackage{xcolor}
\usepackage{setspace}
\usepackage{graphicx}
\usepackage{float}
\usepackage{caption}
\usepackage[colorlinks=true,
            linkcolor=blue,
            citecolor=blue,
            urlcolor=blue]{hyperref}

\definecolor{keywordcolor}{rgb}{0,0,0.6}
\definecolor{commentcolor}{rgb}{0,0.5,0}
\definecolor{stringcolor}{rgb}{0.6,0,0}
\definecolor{backgroundcolor}{rgb}{0.95,0.95,0.95}

\lstset{
    language=Matlab,
    backgroundcolor=\color{backgroundcolor},
    basicstyle=\ttfamily\small,
    keywordstyle=\color{keywordcolor}\bfseries,
    commentstyle=\color{commentcolor}\itshape,
    stringstyle=\color{stringcolor},
    numbers=left,
    numberstyle=\tiny\color{gray},
    stepnumber=1,
    numbersep=8pt,
    frame=none,
    rulecolor=\color{black},
    breaklines=true,
    showstringspaces=false
}

\setstretch{1.2}
\setlength{\parindent}{0pt}
\setlength{\parskip}{0.2\baselineskip}
\newcommand{\MATLAB}{\textsc{Matlab}}

\geometry{margin=2.5cm}
\onehalfspacing

\lstset{
  basicstyle=\ttfamily\small,
  frame=single,
  breaklines=true
}

\setlength{\parindent}{2em}

\begin{document}

\begin{titlepage}
    \begin{center}
		\Huge{Universidade de São Paulo}\\
		\large{Escola de Engenharia de São Carlos}\\ 
        %\vspace{.5cm}
        \large{SEM0530 - Problemas de Engenharia Mecatrônica II}\\ 
        \vspace{5cm}
		\Large{Análise numérica da cinemática do movimento circular de uma partícula}\\ 
		\vspace{0.3cm}
    
    \end{center}
\vspace{5cm}
	\begin{flushleft}
		
		Aluno: Enzo Larger Manfio  Nº USP: 16863802\\
        Docente: Prof. Dr. Kayc Wayhs Lopes\\
	
    \end{flushleft}
	
	\vfill
	\begin{center}
		\vspace{5cm}
			 São Carlos, Maio de 2026
	\end{center}
\end{titlepage}

\newpage
\tableofcontents


\newpage
\section{Introdução}
\par A Integração de funções matemáticas é uma operação fundamental para qualquer aplicação moderna. Contudo, frequentemente é mais generalizável e eficiente calcular uma integral numericamente, ao invés de determina-la analiticamente.
\par Neste trabalho, a integração numérica será utilizada para o estudo da cinemática de uma partícula em movimento circular. Mais especificamente, sabendo a expressão da aceleração tangencial em função da distância percorrida, será desenvolvido um programa na linguagem Matlab que determine os gráficos do módulo da velocidade e aceleração total em função da distância percorrida, bem como determinar o tempo necessário para uma determinada distância percorrida, sendo utilizada a funcionalidade de integração numérica oferecida pela linguagem.
\vspace{0.5\baselineskip}

\section{Metodologia}

\vspace{0.5\baselineskip}
\subsection{Modelagem matemática}
\noindent A aceleração tangencial $a_t$ em função do deslocamento s é dada por \eqref{aceleração_tangencial}:
\begin{equation}
    \label{aceleração_tangencial}
    a_t = 4+0.01Ns - 0.01s^2 + e^{-s/N}.
\end{equation}
Em que N = 2. Com isso, sabendo as velocidade e posição inicial e utilizando a relação diferencial \eqref{matematicos_sofrendo}, é possível expressar o módulo da velocidade tangencial em função do deslocamento:
\begin{align}
    \label{matematicos_sofrendo}
    vdv &= a_tds \\ \nonumber
    \implies \int_{v_0}^{v} vdv &= \int_{s_0}^{s} a_tds \\ \nonumber
    \implies \left. \frac{v^2}{2} \right|_{v_0}^{v} &= \int_{s_0}^{s} a_tds \\
    \label{torriceli_melhorado}
    \implies v &= \sqrt{v_0^2 + 2\int_{s_0}^{s} a_tds}
\end{align}

Com a velocidade inicial sendo dada por \eqref{v_0}.
\begin{equation}
    \label{v_0}
    v_0 = 20 + 0.35\sqrt{N} \text{ }m/s
\end{equation}

\par O vetor aceleração da partícula é composto pela soma dos vetores aceleração tangencial ($a_t$) e radial ($a_r$), os quais são ortogonais entre si, de modo que o módulo da aceleração total, a, é dada por \eqref{pitagoras}:
\begin{equation}
    \label{pitagoras}
    a = \sqrt{a_t^2 + a_r^2}
\end{equation}
$a_t$ é dada por \eqref{aceleração_tangencial}, enquanto $a_r$ é dada por \eqref{aceleração_radial}.
\begin{equation}
    \label{aceleração_radial}
    a_r = \frac{v^2}{r}
\end{equation}
Em que r é o raio da trajetória da partícula, o qual vale 100. Por fim, o tempo t necessário para percorrer uma distância $s - s_0$ pode ser determinado partindo da definição de velocidade \eqref{definição_velocidade}.
\begin{align}
    \label{definição_velocidade}
    v&= \frac{ds}{dt} \\ \nonumber
    \implies \frac{dt}{ds} &= \frac{1}{v} \\ \nonumber
    \implies \int_{s_0}^{s} \frac{dt}{ds}ds &= \int_{v_0}^{v} \frac{1}{v}ds \\ \nonumber
    \implies \int_{0}^{t}dt &= \int_{s_0}^{s} \frac{1}{v}ds \\
    \label{um nome}
    \implies t &= \int_{s_0}^{s} \frac{1}{v}ds
\end{align}
\vspace{0.5\baselineskip}
\subsection{Implementação no \MATLAB}
Em um arquivo próprio, foi definida a função da aceleração tangencial, em acordo com \eqref{aceleração_tangencial}. Com isso, em um outro arquivo próprio foi definida uma função para a velocidade que recebe como argumentos uma determinada função para a aceleração (em função do deslocamento), sua velocidade inicial e seus pontos inicial e final, utilizando \eqref{matematicos_sofrendo}.
\begin{lstlisting}[language=Matlab, 
caption = {Função para determinação da velocidade},
label = {lst:funcao_v_s}]
function v = v_s(s0, v0, at_s, s)
    %Para receber s e retornar v como arrays
    integral_at_ds = arrayfun(@(limite) integral(at_s, s0, limite), s);

    v = sqrt(v0^2 + 2*integral_at_ds);
end
\end{lstlisting}
\par Em um outro arquivo, foi declarada uma função para calcular a aceleração total, recebendo como argumento funções para a velocidade e aceleração tangencial (em função do deslocamento), e o ponto final do movimento, utilizando \eqref{definição_velocidade}.
\begin{lstlisting}[language=Matlab, 
caption = {Função para determinação da aceleração total},
label = {lst:funcao_a_s}]
function a = a_s(v_s, at_s, r, s)
    ar = v_s(s)^2 / r;
    a = sqrt(ar^2 + at_s(s)^2);
end
\end{lstlisting}
\par Por fim, em um outro arquivo foi definida uma função para, recebendo como argumentos uma função para v (em função do deslocamento) e os ponto inicial e final do movimento da partícula, determinar o tempo necessário para realizar este movimento, utilizando \eqref{um nome}.
\begin{lstlisting}[language=Matlab, 
caption = {Função para determinação do tempo},
label = {lst:funcao_a_s}]
function t = t_s(v_s, s0, sf)
    inversa_v = @(s) 1./v_s(s);
    t = integral(inversa_v, s0, sf);
end
\end{lstlisting}
\par Nas três funções, foi utilizada a função \textit{integral}, oferecida pelo Matlab, para calcular as integrais definidas. Ela utiliza o método da Quadratura Adaptativa Global, que divide a integral em subintervalos conforme o erro dos métodos numéricos de integração (trapézio, simpson, etc) aplicados, assim utilizando mais iterações para regiões com maior variação do integrante.
\vspace{0.5\baselineskip}
\section{Resultados}
\subsection{Velocidade}
O gráfico da velocidade em função do deslocamento pode ser visto em \ref{grafico_velocidade}, estando destacado a velocidade da partícula para s = 28,5 m, a qual vale 22.6707 m/s².
\begin{figure}[H]
    \centering
    \includegraphics[width=0.7 \textwidth]{velocidade_deslocamento.pdf}
    \caption{Velocidade em função do deslocamento}
    \label{grafico_velocidade}
\end{figure}
\subsection{Aceleração}
O gráfico da aceleração total em função do deslocamento pode ser visto em \ref{grafico_aceleração}, estando destacado a velocidade da partícula para s = 28,5 m, a qual vale 22.6707 m/s².
\begin{figure}[H]
    \centering
    \includegraphics[width=0.7 \textwidth]{Aceleração_deslocamento.pdf}
    \caption{Aceleração total em função do deslocamento}
    \label{grafico_aceleração}
\end{figure}
\subsection{Tempo}
O tempo necessário para a partícula percorrer 21.5 metros foi calculado como sendo 0.9672 segundos.

\section{Conclusão}
Neste trabalho, foi possível evidenciar o uso de métodos numéricos de integração, com a aplicação destes ao estudo da cinemática do movimento circular de uma partícula, de modo a, a partir de seus parâmetros de velocidade e posição iniciais e função para a aceleração tangencial em função do deslocamento, determinar os gráficos das velocidade e aceleração tangencial em função do deslocamento, bem como determinar o tempo necessário para a partícula realizar um deslocamento arbitrário.

\end{document}